% The CUDA compilation trajectory of nvcc, simplified (Lesson 10).  A
% .cu file is preprocessed twice.  Device path: cicc (front end,
% optimizer, PTX code generation) -> .ptx -> ptxas -> .cubin (SASS),
% repeated per architecture; fatbinary packs PTX and cubin into a fat
% binary.  Host path: cudafe++ separates host code, the host C++ compiler
% compiles it with the fat binary embedded; the linker builds the
% executable.  The device linker nvlink (separate compilation) is omitted.
% Redraws the outline figures m10l1_compileflowgpu and
% m10l1_compileflowgpunvidia (simplified; not NVIDIA's figure).
% Usage: % The CUDA compilation trajectory of nvcc, simplified (Lesson 10).  A
% .cu file is preprocessed twice.  Device path: cicc (front end,
% optimizer, PTX code generation) -> .ptx -> ptxas -> .cubin (SASS),
% repeated per architecture; fatbinary packs PTX and cubin into a fat
% binary.  Host path: cudafe++ separates host code, the host C++ compiler
% compiles it with the fat binary embedded; the linker builds the
% executable.  The device linker nvlink (separate compilation) is omitted.
% Redraws the outline figures m10l1_compileflowgpu and
% m10l1_compileflowgpunvidia (simplified; not NVIDIA's figure).
% Usage: % The CUDA compilation trajectory of nvcc, simplified (Lesson 10).  A
% .cu file is preprocessed twice.  Device path: cicc (front end,
% optimizer, PTX code generation) -> .ptx -> ptxas -> .cubin (SASS),
% repeated per architecture; fatbinary packs PTX and cubin into a fat
% binary.  Host path: cudafe++ separates host code, the host C++ compiler
% compiles it with the fat binary embedded; the linker builds the
% executable.  The device linker nvlink (separate compilation) is omitted.
% Redraws the outline figures m10l1_compileflowgpu and
% m10l1_compileflowgpunvidia (simplified; not NVIDIA's figure).
% Usage: % The CUDA compilation trajectory of nvcc, simplified (Lesson 10).  A
% .cu file is preprocessed twice.  Device path: cicc (front end,
% optimizer, PTX code generation) -> .ptx -> ptxas -> .cubin (SASS),
% repeated per architecture; fatbinary packs PTX and cubin into a fat
% binary.  Host path: cudafe++ separates host code, the host C++ compiler
% compiles it with the fat binary embedded; the linker builds the
% executable.  The device linker nvlink (separate compilation) is omitted.
% Redraws the outline figures m10l1_compileflowgpu and
% m10l1_compileflowgpunvidia (simplified; not NVIDIA's figure).
% Usage: \input{figures/l10-cuda-flow}
% Width: body (120 mm); the drawing is about 116 mm wide.
\begin{tikzpicture}[x=1mm, y=1mm, font=\scriptsize\sffamily,
    >={Stealth[length=1.5mm]},
    tool/.style={draw, fill=ln-accent!20, minimum height=8mm, align=center,
                 inner sep=1.5pt},
    file/.style={draw, fill=ln-box, rounded corners=1.5pt, minimum height=6mm,
                 align=center, inner sep=1.5pt, font=\scriptsize\ttfamily}]
  % source
  \node[file, minimum height=10mm] (cu) at (6,20) {x.cu};
  % device row
  \node[tool] (dpp) at (22,36) {\iflangja{プリ\\プロセッサ}{pre-\\processor}};
  \node[tool, fill=ln-source!25] (cicc) at (44,36) {cicc\\[-1pt]
    {\scriptsize\sffamily\iflangja{FE＋最適化＋PTX生成}{FE + opt.\ + PTX gen.}}};
  \node[file] (ptx) at (65,36) {x.ptx};
  \node[tool, fill=ln-source!25] (ptxas) at (81,36) {ptxas};
  \node[file] (cubin) at (100,36) {x.cubin\\[-1pt]{\scriptsize\sffamily (SASS)}};
  % fatbinary
  \node[tool] (fatb) at (100,20) {fatbinary};
  % host row
  \node[tool] (hpp) at (22,4) {\iflangja{プリ\\プロセッサ}{pre-\\processor}};
  \node[tool] (fe) at (44,4) {cudafe++};
  \node[tool] (cxx) at (71,4) {\iflangja{ホスト C++\\コンパイラ}{host C++\\compiler}};
  \node[file] (obj) at (91,4) {x.o};
  \node[tool] (ld) at (108,4) {\iflangja{リンカ}{linker}};
  \draw[->] (cu.north) |- (dpp.west);
  \draw[->] (cu.south) |- (hpp.west);
  \draw[->] (dpp) -- (cicc);
  \draw[->] (cicc) -- (ptx);
  \draw[->] (ptx) -- (ptxas);
  \draw[->] (ptxas) -- (cubin);
  \draw[->] (cubin) -- (fatb);
  \draw[->] (ptx.south) |- (fatb.west);
  \draw[->] (hpp) -- (fe);
  \draw[->] (fe) -- (cxx);
  \draw[->] (cxx) -- (obj);
  \draw[->] (obj) -- (ld);
  \draw[->, dashed] (fatb.south) -- ++(0,-4) -| node[pos=0.25, above, inner sep=1pt]
    {\iflangja{ファットバイナリを埋め込む}{fat binary embedded}} (cxx.north);
  \draw[->] (ld.south) -- ++(0,-6) node[below, inner sep=1pt]
    {\iflangja{実行ファイル}{executable}};
  % lane labels
  \node[anchor=south west, text=ln-muted] at (13,41.5)
    {\iflangja{デバイスコード（アーキテクチャごとに繰り返す）}{device code (repeated per architecture)}};
  \node[anchor=north west, text=ln-muted] at (13,-1.5)
    {\iflangja{ホストコード}{host code}};
\end{tikzpicture}

% Width: body (120 mm); the drawing is about 116 mm wide.
\begin{tikzpicture}[x=1mm, y=1mm, font=\scriptsize\sffamily,
    >={Stealth[length=1.5mm]},
    tool/.style={draw, fill=ln-accent!20, minimum height=8mm, align=center,
                 inner sep=1.5pt},
    file/.style={draw, fill=ln-box, rounded corners=1.5pt, minimum height=6mm,
                 align=center, inner sep=1.5pt, font=\scriptsize\ttfamily}]
  % source
  \node[file, minimum height=10mm] (cu) at (6,20) {x.cu};
  % device row
  \node[tool] (dpp) at (22,36) {\iflangja{プリ\\プロセッサ}{pre-\\processor}};
  \node[tool, fill=ln-source!25] (cicc) at (44,36) {cicc\\[-1pt]
    {\scriptsize\sffamily\iflangja{FE＋最適化＋PTX生成}{FE + opt.\ + PTX gen.}}};
  \node[file] (ptx) at (65,36) {x.ptx};
  \node[tool, fill=ln-source!25] (ptxas) at (81,36) {ptxas};
  \node[file] (cubin) at (100,36) {x.cubin\\[-1pt]{\scriptsize\sffamily (SASS)}};
  % fatbinary
  \node[tool] (fatb) at (100,20) {fatbinary};
  % host row
  \node[tool] (hpp) at (22,4) {\iflangja{プリ\\プロセッサ}{pre-\\processor}};
  \node[tool] (fe) at (44,4) {cudafe++};
  \node[tool] (cxx) at (71,4) {\iflangja{ホスト C++\\コンパイラ}{host C++\\compiler}};
  \node[file] (obj) at (91,4) {x.o};
  \node[tool] (ld) at (108,4) {\iflangja{リンカ}{linker}};
  \draw[->] (cu.north) |- (dpp.west);
  \draw[->] (cu.south) |- (hpp.west);
  \draw[->] (dpp) -- (cicc);
  \draw[->] (cicc) -- (ptx);
  \draw[->] (ptx) -- (ptxas);
  \draw[->] (ptxas) -- (cubin);
  \draw[->] (cubin) -- (fatb);
  \draw[->] (ptx.south) |- (fatb.west);
  \draw[->] (hpp) -- (fe);
  \draw[->] (fe) -- (cxx);
  \draw[->] (cxx) -- (obj);
  \draw[->] (obj) -- (ld);
  \draw[->, dashed] (fatb.south) -- ++(0,-4) -| node[pos=0.25, above, inner sep=1pt]
    {\iflangja{ファットバイナリを埋め込む}{fat binary embedded}} (cxx.north);
  \draw[->] (ld.south) -- ++(0,-6) node[below, inner sep=1pt]
    {\iflangja{実行ファイル}{executable}};
  % lane labels
  \node[anchor=south west, text=ln-muted] at (13,41.5)
    {\iflangja{デバイスコード（アーキテクチャごとに繰り返す）}{device code (repeated per architecture)}};
  \node[anchor=north west, text=ln-muted] at (13,-1.5)
    {\iflangja{ホストコード}{host code}};
\end{tikzpicture}

% Width: body (120 mm); the drawing is about 116 mm wide.
\begin{tikzpicture}[x=1mm, y=1mm, font=\scriptsize\sffamily,
    >={Stealth[length=1.5mm]},
    tool/.style={draw, fill=ln-accent!20, minimum height=8mm, align=center,
                 inner sep=1.5pt},
    file/.style={draw, fill=ln-box, rounded corners=1.5pt, minimum height=6mm,
                 align=center, inner sep=1.5pt, font=\scriptsize\ttfamily}]
  % source
  \node[file, minimum height=10mm] (cu) at (6,20) {x.cu};
  % device row
  \node[tool] (dpp) at (22,36) {\iflangja{プリ\\プロセッサ}{pre-\\processor}};
  \node[tool, fill=ln-source!25] (cicc) at (44,36) {cicc\\[-1pt]
    {\scriptsize\sffamily\iflangja{FE＋最適化＋PTX生成}{FE + opt.\ + PTX gen.}}};
  \node[file] (ptx) at (65,36) {x.ptx};
  \node[tool, fill=ln-source!25] (ptxas) at (81,36) {ptxas};
  \node[file] (cubin) at (100,36) {x.cubin\\[-1pt]{\scriptsize\sffamily (SASS)}};
  % fatbinary
  \node[tool] (fatb) at (100,20) {fatbinary};
  % host row
  \node[tool] (hpp) at (22,4) {\iflangja{プリ\\プロセッサ}{pre-\\processor}};
  \node[tool] (fe) at (44,4) {cudafe++};
  \node[tool] (cxx) at (71,4) {\iflangja{ホスト C++\\コンパイラ}{host C++\\compiler}};
  \node[file] (obj) at (91,4) {x.o};
  \node[tool] (ld) at (108,4) {\iflangja{リンカ}{linker}};
  \draw[->] (cu.north) |- (dpp.west);
  \draw[->] (cu.south) |- (hpp.west);
  \draw[->] (dpp) -- (cicc);
  \draw[->] (cicc) -- (ptx);
  \draw[->] (ptx) -- (ptxas);
  \draw[->] (ptxas) -- (cubin);
  \draw[->] (cubin) -- (fatb);
  \draw[->] (ptx.south) |- (fatb.west);
  \draw[->] (hpp) -- (fe);
  \draw[->] (fe) -- (cxx);
  \draw[->] (cxx) -- (obj);
  \draw[->] (obj) -- (ld);
  \draw[->, dashed] (fatb.south) -- ++(0,-4) -| node[pos=0.25, above, inner sep=1pt]
    {\iflangja{ファットバイナリを埋め込む}{fat binary embedded}} (cxx.north);
  \draw[->] (ld.south) -- ++(0,-6) node[below, inner sep=1pt]
    {\iflangja{実行ファイル}{executable}};
  % lane labels
  \node[anchor=south west, text=ln-muted] at (13,41.5)
    {\iflangja{デバイスコード（アーキテクチャごとに繰り返す）}{device code (repeated per architecture)}};
  \node[anchor=north west, text=ln-muted] at (13,-1.5)
    {\iflangja{ホストコード}{host code}};
\end{tikzpicture}

% Width: body (120 mm); the drawing is about 116 mm wide.
\begin{tikzpicture}[x=1mm, y=1mm, font=\scriptsize\sffamily,
    >={Stealth[length=1.5mm]},
    tool/.style={draw, fill=ln-accent!20, minimum height=8mm, align=center,
                 inner sep=1.5pt},
    file/.style={draw, fill=ln-box, rounded corners=1.5pt, minimum height=6mm,
                 align=center, inner sep=1.5pt, font=\scriptsize\ttfamily}]
  % source
  \node[file, minimum height=10mm] (cu) at (6,20) {x.cu};
  % device row
  \node[tool] (dpp) at (22,36) {\iflangja{プリ\\プロセッサ}{pre-\\processor}};
  \node[tool, fill=ln-source!25] (cicc) at (44,36) {cicc\\[-1pt]
    {\scriptsize\sffamily\iflangja{FE＋最適化＋PTX生成}{FE + opt.\ + PTX gen.}}};
  \node[file] (ptx) at (65,36) {x.ptx};
  \node[tool, fill=ln-source!25] (ptxas) at (81,36) {ptxas};
  \node[file] (cubin) at (100,36) {x.cubin\\[-1pt]{\scriptsize\sffamily (SASS)}};
  % fatbinary
  \node[tool] (fatb) at (100,20) {fatbinary};
  % host row
  \node[tool] (hpp) at (22,4) {\iflangja{プリ\\プロセッサ}{pre-\\processor}};
  \node[tool] (fe) at (44,4) {cudafe++};
  \node[tool] (cxx) at (71,4) {\iflangja{ホスト C++\\コンパイラ}{host C++\\compiler}};
  \node[file] (obj) at (91,4) {x.o};
  \node[tool] (ld) at (108,4) {\iflangja{リンカ}{linker}};
  \draw[->] (cu.north) |- (dpp.west);
  \draw[->] (cu.south) |- (hpp.west);
  \draw[->] (dpp) -- (cicc);
  \draw[->] (cicc) -- (ptx);
  \draw[->] (ptx) -- (ptxas);
  \draw[->] (ptxas) -- (cubin);
  \draw[->] (cubin) -- (fatb);
  \draw[->] (ptx.south) |- (fatb.west);
  \draw[->] (hpp) -- (fe);
  \draw[->] (fe) -- (cxx);
  \draw[->] (cxx) -- (obj);
  \draw[->] (obj) -- (ld);
  \draw[->, dashed] (fatb.south) -- ++(0,-4) -| node[pos=0.25, above, inner sep=1pt]
    {\iflangja{ファットバイナリを埋め込む}{fat binary embedded}} (cxx.north);
  \draw[->] (ld.south) -- ++(0,-6) node[below, inner sep=1pt]
    {\iflangja{実行ファイル}{executable}};
  % lane labels
  \node[anchor=south west, text=ln-muted] at (13,41.5)
    {\iflangja{デバイスコード（アーキテクチャごとに繰り返す）}{device code (repeated per architecture)}};
  \node[anchor=north west, text=ln-muted] at (13,-1.5)
    {\iflangja{ホストコード}{host code}};
\end{tikzpicture}
